\section{Understanding the Emergent Language}
\label{sec:understanding-emergent-language}

%\textbf{AL: I'm now wondering if we should have a 3-section piece, with understanding/measuring being its own thing after the new simulations and the agent coordination?}
With more realistic simulations, understanding
what is going on becomes more difficult. Even when we are confident
that genuine communication is taking place (Section
\ref{sec:measuring-communication}), it is difficult to decode messages
by simple inspection. We do not know if and how the messages produced
by the agents should be segmented into ``words''. We might only have
vague conjectures about what they refer to. If there are multiple
turns, we do not know which turns are mostly information exchanges,
and which, if any, absolve other pragmatic functions (e.g., asking for more information). We cannot even
trust the agents to use symbols in a consistent way across contexts
and turns \cite{Bogin:etal:2018}. The enterprise is akin to
linguistic fieldwork, except that we are dealing with an alien race,
with no guarantees that universals of human communication will
apply.% There has thus been extensive work focusing
% on decoding the agents' emergent language.

Indeed, in spite of task success, the emerging language can have
counter-intuitive properties. \citeA{Kottur:etal:2017} considered
agents playing a referential game in which they must communicate about
object attributes and values (e.g, \emph{color: blue}, \emph{shape:
  round}). The agents have difficulties converging to the intuitive
coding scheme in which distinct symbols unambiguously denote single
attributes or values (i.e., a word for \emph{color}, a word for
\emph{blue}, etc). Such code will only emerge when the set of
available symbols is greatly limited and the memory of one of the
agents is ablated, pointing to memory bottlenecks as a possible
bias to be injected into deep networks for more natural languages to
emerge \cite<see also>{Resnick:etal:2020}. \citeA{Bouchacourt:Baroni:2018} replicated
the game of \citeA{Lazaridou:etal:2017} with the surprising results
illustrated in Fig.~\ref{fig:gaussian}.
% considered a multi-turn game based on
% single-symbol exchanges. An agent sees an object, and the other is
% tasked with reporting the values of two attributes of the object
% (e.g., color and shape). If the agents are allowed to use a large set
% of symbols, they will learn to associate symbols to whole attribute
% value combinations (e.g., a single symbol meaning
% \emph{blue-square}). This strategy results in a language that cannot
% generalize to new combinations of attributes or values. More
% interestingly, even when the available symbol cardinality is reduced,
% the agents do not converge to the intuitive strategy of using them to
% denote single attributes and values. Instead, the agent that needs to
% ask about attributes saves on vocabulary usage by letting the same
% symbol refer to attribute combinations in an order-insensitive way
% (e.g., one symbol denotes both \emph{color-shape} and
% \emph{shape-color}). The other agent uses symbol sequences to refer to
% different value combinations in an arbitrary, context-sensitive way
% (e.g., symbol \emph{A} followed by \emph{B} might refer to
% \emph{blue}, but when followed by \emph{C} it might refer to
% \emph{square}). Kottur et al.~finally show that, only when the memory
% of one agent is removed, so that the latter cannot rely on context,
% the emergent language words consistently denote single attributes and
% values. This comes, however, at the cost of agent lobotomy. Bogin et
% al.~\cite{Bogin:etal:2018} also found that, with standard training
% methods and a sufficiently complex environment, agents will not
% develop a consistent language where symbols have the same referent
% independently of context. 

% Figure Legends (250 words per legend)
% - Should always have a short, explanatory title (as well as legend).
% - Legends must explain the figure fully without reference to the text.
% - If applicable, the original source of previously published material must be acknowledged in the Figure legend.
% - Figures must be cited in the main text (box figures must be cited in the box). 

\begin{figure}[p]
  \centering
\includegraphics[width=7cm]{figures/gaussian}
\caption{\textbf{Training and test inputs in the referential game of
    \cite{Bouchacourt:Baroni:2018}.} Two agents were trained to play
  the  game of \citeA{Lazaridou:etal:2017} (see
  Fig.~\ref{fig:referential-game}). During training, the agents were
  exposed to the same data as in the original study, that is, pairs of
  pictures of instances of about 500 distinct objects (top row). At
  test time, however, the agents were made to play the game with blobs
  of Gaussian noise (bottom row). They were able to
  communicate about them nearly as well as about the training
  pictures. This shows that the language emerging in this game does
  not involve ``words'' referring to generic concepts, but rather
  \emph{ad-hoc} signals, probably carrying comparative information
  about shallow visual properties of the images. Bottom row reproduced
  from \citeA{Bouchacourt:Baroni:2018} by permission.}
  \label{fig:gaussian}
\end{figure}


Essentially, agents will develop a code that is sufficient to solve
the task at hand, and hoping that such code will possess further
desirable characteristics is wishful thinking. Consider the task in
the original game of \citeA{Lazaridou:etal:2017}. The agents must
discriminate pairs of pictures depicting instances of 500
categories. The agents could achieve this by developing human-like
names for the categories, but a low-level strategy relying on, say,
comparing average pixel intensity in patches of the two images might
require as few symbols as 2. In this respect, the
agents' language is, paradoxically, ``too human'', in the sense that
it evolved to minimize effort, while remaining adequate for the task
at hand \cite{Gibson:etal:2019}. % ; compare to the tendency of humans in
% restricted communication setups to converge to \emph{ad-hoc} naming
% conventions, or ``conceptual pacts'',
% \cite{Brennan:Clark:1996}). 
Indeed, \citeA{Kharitonov:etal:2020} showed
that the way deep agent emergent languages partition their
meaning space displays the same tendency towards complexity minimization
that is pervasive in human language.

%Still on the topic of effort minimization,
\citeA{Chaabouni:etal:2019} studied whether agent language
% In part because of the aforementioned difficulties, other studies
% focus on arguably even more basic properties of natural language than
% compositionality. For example, it is well known that human languages
% and other animal communication systems
exhibit an inverse correlation between word frequency and word length,
so that the signals that need to be used more often are also the
shortest, as universally found in natural languages
\cite{Zipf:1949,Strauss:etal:2007,FerrerICancho:etal:2013}. They discovered that deep agents trained with a referential game where inputs
have a skewed distribution similar to natural language actually
develop a significantly \emph{anti-efficient} code, in which the
most frequent inputs are associated to the longest messages. The
effect is explained by the lack of an articulatory effort
minimization bias in networks, that are thus only subject to a
``perceptual'' pressure favoring longer messages, as they are easier to discriminate.
% would also develop a language exhibiting this
% property. Not only this was not the case: the agents developed a
% positively \emph{anti-efficient} code, in which the most frequent
% inputs were associated to the longest messages. They showed that this
% surprising pattern emerges because listening agents have an innate
% preference for longer messages, as they are more discriminative. The
% same might be true of human and animal perceivers. However, in
% biological agents this is counterbalanced by the speaker preference
% for conciseness, as message articulation is costly. Since artificial
% networks are not subject to comparable metabolic constraints, there is
% no similar contrasting pressure, and  codes that only favour the listener can
% easily come about. Paradoxically, by providing a counterexample to its universal
% emergence, this finding supports the view that the inverse
% frequency/length correlation attested in humans and animal
% communication systems is indeed due to efficiency factors. If it was a
% trivial effect of random string assignment processes, as sometimes
% conjectured (\cite{Miller:1957,delPradoMartin:2013}), then it should
% arise in any successful communication system dealing with a skewed
% input referent distribution.

% As another example, Chaabouni et al.~\cite{Chaabouni:etal:2019} studied whether deep agents have innate
% word order biases. The agents were trained to use simple languages
%  displaying or lacking certain properties. Relative ease of learning by
% individual agents, as well as language drift in cross-generation
% transmission, were used to probe how natural the relevant
% characteristics were for the agents. Chaabouni and colleagues found
% deep agents to have priors in line with some universal preferences of
% human language, such as a tendency to reduce word order variation, and
% avoidance of long-distance dependencies. Intriguingly, agents, like
% humans, prefer symbols referring to a sequence of events to be encoded
% iconically, but in \emph{reverse} order (that is, saying something
% like \emph{turn left, go straight} to mean ``first go straight, then
% turn left''). Finally, in another example of how artificial agents are
% not subject to the same least-effort pressures that shape human
% language, they found that the agents did not have a tendency to shed
% redundant cues of the same information (such as double role marking
% by means of word order and overt suffixes). All in all, the results of this study further confirm the mixed picture about the relation between deep agent codes (and the underlying biases that lead to their emergence) and natural language.  

% box communication
\subsection{Measuring the Degree of Effective Communication}
\label{sec:measuring-communication}
  As simulations move beyond referential games (where task success
  trivially depends on establishing a communication code), to complex
  environments where communication plays an auxiliary function (e.g.,
  \citeA{Das:etal:2019}; Fig.~\ref{fig:environments}(d)), the first
   question to ask when analyzing an emergent language is whether
  it is actually been used in any meaningful way by the agents.  As
  clearly discussed by \citeA{Lowe:etal:2019}, just ablating the
  language channel and showing a drop in task success does not prove
  much, as the extra capacity afforded by the channel architecture
  might have helped the agents' learning process without being used to
  establish communication. The same paper proposes a classification of
  measures to detect the presence of genuine
  communication. \textit{Positive signaling} captures the extent to
  which information about the sender states, observations and actions
  are expressed in its signals. \textit{Positive listening} captures
  the extent to which a signal impacts the receiver's states and
  behaviour. Examples of positive signaling include \textit{context
    independence}~\cite{Bogin:etal:2018} and \textit{speaker
    consistency}~\cite{Jaques:etal:2019}. The former measures the
  degree of alignment between messages and task-related concepts,
  whereas the latter measures, through mutual information, the
  alignment between an agent messages and its actions. Positive
  signaling gives no guarantee of communication, since the receiver could
  be ignoring the sender messages, no matter how informative they might
  be. An example of positive listening is the \textit{instantaneous
    coordination} measure of~\citeA{Jaques:etal:2019}, which uses
  mutual information to quantify correlation between 
  sender messages and receiver actions. Instead,~\citeA{Lowe:etal:2019}
  propose to use \emph{causal influence of communication}, a
  quantity that measures the causal relationship between sender
  messages and receiver actions.   The authors show that only a high causal influence
  of communication is both necessary and sufficient for
  positive listening, and thus communication.

  The call for caution
  is not just hypothetical. Several studies have reported how agents
  can easily converge to non-verbal or degenerate strategies, even
  when it would seem that communication is taking place. For example,
   agents might learn to exchange information simply through the
  number of turns they take before ending the game, irrespective of
  what they actually say
  \cite{Cao:etal:2018,Bouchacourt:Baroni:2019}.

\subsection{Compositionality}

% Starting from similar considerations, much work on analyzing emergent
% languages has focused on which environmental aspects and biases
% (``innate'' or injected into the agent architectures) force the
% development of a code possessing some properties of interest.

Much analytical work in the area has focused on compositionality, as the latter is seen both as a
fundamental feature of natural language whose evolutionary origins are
unclear \cite{Bickerton:2014,Townsend:etal:2018}, and as a pre-condition
for an emergent language to generalize at scale.
% There
% has been in particular widespread interest in the conditions under
% which the emergent language is compositional. This interest is
% theoretically motivated by the idea that compositionality is a
% fundamental, perhaps unique feature of human language, but its
% evolutionary origins are not clear
% (\cite{Bickerton:2014,Townsend:etal:2018}). From an applied angle,
% compositionality might be a pre-condition for a language to be able
% to generalize at scale.

% In linguistics, compositionality is traditionally defined as the
% property by which the meaning of a composite linguistic expression is
% a function of the meanings of its parts (and systematic meaning
% composition rules) (\cite{Partee:2004}). In the context of current
% language emergence work, composite meanings are typically restricted
% to simple conjunctions of attribute values (is the message denoting
% \emph{red triangles} predictable from messages denoting other
% \emph{red} things and messages denoting \emph{triangles} of other
% colors?). Even so, it turns out to be very hard to assess how
% compositional an emergent language is. First, we have not yet
% developed systematic techniques to segment messages into component
% parts (note that there is no guarantee that such parts should be
% combined by agents concatenatively). Second, a black-or-white view of compositionality, such as the
% one suggested by the standard linguists' definition, is too coarse. It
% is likely that different emergent languages will show different
% degrees of compositionality, and missing this gradient might imply missing
% important insights onto what makes a language take the first
% steps towards full compositionality.

The simplest way to probe for compositionality in an emergent protocol
is to test whether agents can use it to denote novel composite
meanings, e.g., can they refer to \emph{blue squares} on first encounter, if
they have seen other \emph{blue} and \emph{square} things during
training \cite{Choi:etal:2018}. This assumes that a
compositional encoding is necessary to generalize. However, a few recent papers have intriguingly reported %
% address these challenges is to directly test the
% ability of agents to refer to novel composite meanings as a
%  proxy of compositionality (e.g., can the agents
% successfully refer to a \emph{blue square} the first time they are
% exposed to it, if in the training phase they encountered other
% \emph{blue} and \emph{square} things?)
% (e.g., \cite{Choi:etal:2018}). However, this approach is assuming that
% something akin to the linguists' notion of compositionality is
% necessary for generalization, which is actually the sort of hypothesis
% we would like to \emph{validate} with our simulations. Indeed, one of
% the most intriguing observations coming from the field is
that emergent languages can support generalization to novel composite meanings \emph{without} conforming to even weak notions of
compositionality \cite{Lazaridou:etal:2018,Andreas:2019,Chaabouni:etal:2020}. %  an
% empirical result that we hope to be explored in much more detail in
% the future. Moreover, generalization might tell us at best that some
% form of compositionality is in place, but it won't give us any cue on
% the kind of compositional processes that a language is making use of.
\citeA{Lazaridou:etal:2018} re-introduced to the emergent language
community the \emph{topographic similarity} score from earlier work on
language emergence \cite{Brighton:2006}. Given ways to measure
distances between meanings and between forms, topographic similarity
is the correlation between all possible meaning pair distances and the
distances of the corresponding message pairs. It captures the intuition that compositionality involves a systematic
relation between form and meaning similarity. However, it does not tell us
anything about the nature of the specific compositional processes present
in a language. %
% and it does so by
% providing a continuous score. Moreover, minimum edit distance is
% relatively robust to non-concatenative changes in form, as it can
% capture ways to compose messages involving deletions and
% insertions. Conversely, exactly because of its generality, a high
% topographic similarity does not tell us much about the nature of the
% compositional processes that are present in a language.
\citeA{Andreas:2019}
proposed a method to quantify to what extent an emergent language reflects specific types of compositional structure in its
input. Unfortunately, the method only works if we have a concrete
hypothesis about the underlying composition function, that is, it can only be used to test whether a language conforms to an underlying compositional grammar if we are able to precisely specify this grammar. This limits its practical
applicability. Finally, \citeA{Chaabouni:etal:2020} recently established a link between compositionality and the notion of disentanglement in representation learning \cite{Suter:etal:2019}, and proposed to use methods to quantify disentanglement from that literature in order to measure the degree of compositionality of emergent codes.

Equipped with similar tools, various studies have uncovered different aspects
of compositionality in emergent languages. For example, \citeA{Lazaridou:etal:2018}
found that compositionality more easily emerges when objects are
represented symbolically as sets of attribute-value pairs, than when
they are more realistically represented as synthetic 3D shapes. %
%
% let two agents play a referential game where the speaker was presented
% with a target object and produced a (multi-symbol) message, and the
% receiver had to point to the target amidst distractors. They found
% that when objects were represented as symbolic attribute-value sets,
% the emergent language exhibited at least a certain degree of
% compositional structure, as reflected by topographic similarity. The
% picture was more mixed when realistic perceptual processing was added
% to the simulation, with objects being represented as 3D geometrical
% shapes in synthesized scenes, that were given to the agents as raw
% pixel maps \textbf{(add figure to better explain this?)}. On the one
% hand, some degree of compositionality was also present in at least
% some of the experiments run under this regime. However, the pattern
% was now quite brittle. In particular, reducing the number of
% distractors led to the emergence of non-compositional,
% hard-to-interpret languages, probably focusing on low-level aspects of
% the comparison between the target and the few distractor images
% left. Again, the take-home message might be that intuitive languages
% (in this case, languages characterizing geometric shapes in terms of
% discrete visual attributes) will only emerge if the task is
% challenging enough to justify the complexity of the solution (the more
% compositional language that emerged in the simulations uses more than
% 1000 distinct messages, compared to a dozen or less for the degenerate
% ones). As Lazaridou and colleagues point out, the tendency to converge
% to \emph{ad-hoc} solutions that are good-enough for the task at hand
% is not an unnatural property of deep agents, but rather an instance of
% the same efficiency pressure that leads humans to establish
% specialized naming conventions (``conceptual pacts'') in specific
% communicative settings (\cite{Brennan:Clark:1996}).
% Other studies have explored other aspects of compositionality in
% emergent languages. For example,
% Mordatch and Abbeel
\citeA{Mordatch:Abbeel:2018}  studied
the code emerging in a community of agents moving and
acting in a shared grid-world. Each agent was assigned a goal, that
could involve having another agent moving to a landmark position. An
order-insensitive concatenative language emerged, where agents would
refer to actions, their agent and targets by juxtaposing specialized
symbols (e.g., one message could be: \emph{goto blue-agent
  red-landmark}, or, equivalently, \emph{red-landmark goto
  blue-agent}). % When agents were deprived of the
% communication channel, they resorted to non-verbal strategies, such as
% pointing or pushing, to direct each other.

Despite many intriguing empirical observations, our characterization
of which architectural biases and environmental pressures favour the
emergence of compositionality (or other linguistic properties) is
still very sketchy. A strong result obtained both with humans and in
pre-deep-learning computational simulations is that generational
transmission of language favors compositionality
\cite<e.g.,>{Kirby:Hurford:2002,Kirby:etal:2014}, an observation recently
confirmed for deep agents \cite{Li:Bowling:2019,Ren:etal:2019}.
Moreover, recent results from experiments with humans demonstrate that
larger communities of speakers evolve more systematic languages
\cite{Raviv:etal:2019a,Raviv:etal:2019b}, suggesting the need to
move away from two-partner agent setups, an observation
beginning to find its way into deep agent research
\cite{Graesser:etal:2019,Tieleman:etal:2019}.

Other priors for compositionality that have been proposed and at least partially
empirically validated include input representations
\cite{Lazaridou:etal:2018}, agent and channel capacity
\cite{Kottur:etal:2017,Mordatch:Abbeel:2018,Resnick:etal:2020}, and
specific training strategies, such as letting the agents simulate
other agents' understanding of one's language
\cite{Choi:etal:2018}. We still lack, however, systematic
experiments establishing which of these conditions are necessary,
which are sufficient, and how they interact, possibly along the lines
of earlier systematizing work such as \citeA{Bratman:etal:2010}.
% Overall, though, our characterization of the emergence of compositionality in
% deep agent communication is much murkier than in the small controlled
% experiments from the previous decades. First, a number of possible
% priors favoring compositionality have found preliminary support,
% including, besides generational transmission, constraints on the %environmental pressures
% environment and input representations (\cite{Lazaridou:etal:2018}),
% agent and channel capacity
% (\cite{Kottur:etal:2017,Mordatch:Abbeel:2018}), and specific training
% strategies, such as letting the agents simulate other agents'
% understanding of one's language (\cite{Choi:etal:2018}). We still
% lack systematic experiments establishing which
% of these conditions are necessary, which are sufficient, and how they
% interact. Second, except in the tiniest experiments, we never really
% observe the emergence of a fully compositional language. Rather, emergent languages tend at best to mix messages (or message substrings) that can be
% transparently segmented into components and messages that are
% partially or completely opaque. Further analysis of this intriguing
% result is needed.

